\section{Related Work}
\label{sec:related_work}

\paragraph{Realized volatility and the HAR class.}
Realized variance as a nonparametric estimator of integrated variance from
high-frequency returns is due to \citet{AndersenBollerslevDieboldLabys2003},
and the heterogeneous autoregressive model of \citet{Corsi2009} remains the
reference specification: a regression of future realized variance on
daily, weekly and monthly trailing averages, motivated as a cascade of
market participants operating at different horizons and as a parsimonious
approximation to long memory. Extensions add jump components
\citep{AndersenBollerslevDiebold2007, BarndorffNielsenShephard2004},
correct for measurement error in the realized proxy
\citep{BollerslevPattonQuaedvlieg2016}, or couple returns and realized
measures in a single likelihood \citep{HansenHuangShek2012}.

Our reading of HAR in Section~\ref{sec:har_quadrature} differs from the
cascade interpretation in a way that is empirically consequential. We treat
the lag ladder not as a set of economically distinct horizons but as a
quadrature grid for a single continuous decay kernel, which predicts that
refining the grid should pay monotonically. Table~\ref{tab:ladder} confirms
this on our panel, and the conventional daily/weekly/monthly structure is
\emph{worse} than a denser geometric ladder. To our knowledge the lag grid has
not previously been treated as a tunable quadrature in this literature.

\paragraph{Hawkes processes and endogenous volatility.}
Self-exciting point processes \citep{Hawkes1971, HawkesOakes1974} model
activity as a superposition of exogenous arrivals and endogenously triggered
offspring, with the branching ratio $n = \int \varphi$ measuring the
endogenous share. Applications to financial data are surveyed by
\citet{BacryMastromatteoMuzy2015}. The empirical kernel is consistently found
to be power-law, and estimates of $n$ close to unity have been read as
evidence that markets operate near criticality
\citep{HardimanBercotBouchaud2013}; \citet{FilimonovSornette2012} argue that
such estimates are inflated when the background intensity is time-varying and
unmodelled, which is directly relevant to intraday data with strong
periodicity \citep{AndersenBollerslev1997}.

Two extensions bear on our results. Nonlinear Hawkes processes admit
inhibition \citep{BremaudMassoulie1996}, which is the only class in which a
kernel with negative lobes---such as the one we estimate---is well defined.
Quadratic Hawkes \citep{BlancDonierBouchaud2017} makes the feedback a
quadratic form in the past, generating clustering, fat tails and the Zumbach
effect \citep{Zumbach2009} endogenously; it predicts extractable structure
saturating at pairwise order, which is what Section~\ref{sec:saturation}
measures. For the multivariate case, estimation of low-rank and sparse kernel
matrices is treated by \citet{BacryGaiffasMuzy2015} and
\citet{AchabBacryGaiffas2017}.

\paragraph{Rough volatility.}
\citet{JaissonRosenbaum2016} show that nearly unstable Hawkes processes with
heavy-tailed kernels converge under rescaling to rough fractional processes,
supplying a microstructural foundation for the rough-volatility models of
\citet{GatheralJaissonRosenbaum2018} and \citet{ElEuchRosenbaum2019}; see also
\citet{BennedsenLundePakkanen2022} on separating short- and long-horizon
behaviour. This literature is the closest competitor to our backbone, since it
too predicts a power-law prediction kernel, and Section~\ref{sec:rough} notes
that our battery does not yet contain a fractional-kernel arm. We regard that
as the most consequential missing benchmark in the paper.

\paragraph{Machine learning for volatility forecasting.}
\citet{gu2020empirical} establish the modern protocol for large-scale model
comparison in asset pricing, and \citet{christensen2023machine} apply it to
realized volatility, reporting gains for nonlinear methods. Results in this
area disagree about whether nonlinearity helps, and by how much. Our
Section~\ref{sec:tuning_leak} contributes a measured account of one mechanism
behind that disagreement: a hyperparameter search scored on the evaluation
panel produced, on our own data, an apparent model-class reversal roughly an
order of magnitude larger than any effect we can defend under a causal
protocol. We therefore read cross-study disagreement as partly a protocol
artifact rather than as substantive disagreement about volatility.

Our finding that dense shrinkage dominates sparse selection is consistent with
the weak-and-dense-signal view of return predictability, and our attribution
machinery follows the Frisch--Waugh--Lovell tradition
\citep{FrischWaugh1933, Lovell1963}. For inference we use Diebold--Mariano
tests \citep{DieboldMariano1995} with the small-sample correction of
\citet{HarveyLeybourneNewbold1997}, the model confidence set of
\citet{HansenLundeNason2011}, and QLIKE as the primary loss on the robustness
grounds of \citet{Patton2011}. Section~\ref{sec:calibration} reports a
property of that loss on our panel---that it systematically rewards attenuated
forecasts---which we have not seen quantified elsewhere in this setting.
